% GENERATED FILE. Edit canonical lessons, then run npm run generate:books.
% canonical-source-hash: fnv1a64:3286bf5fd2f356de
% canonical-lessons: GU-C359-ijjat, GU-C359-gaurav, GU-C359-saram, GU-C359-abhiman, GU-C359-namvu

\chapter{Honour and Shame}
\label{ch:gu-honour-and-shame}
\hlchaptermodality{359}

\begin{chapteropening}
\textbf{By the end of this chapter:} \emph{I can talk about honour, pride and shame.}

Complete the last lesson of chapter 359: I can talk about honour, pride and shame.
\end{chapteropening}

\section[\texorpdfstring{\gu{ઇજ્જત} (ijjat)}{ઇજ્જત (ijjat)}]{\gu{ઇજ્જત} (ijjat) — one's good name; honour}

\label{lesson:GU-C359-ijjat}



\begin{quote}
Before the new one: say the Gujarati for cooperation, then the Gujarati for to keep up.
\end{quote}

\subsection*{You'll want to know: \gu{ઇજ્જત}}
\textbf{\gu{ઇજ્જત}} — \emph{ijjat} — \textquotedblleft{}one's good name; honour\textquotedblright{}.

\begin{quote}
\gu{ગામમાં} \gu{તેના} \gu{કુટુંબની} \gu{ઘણી} \gu{ઇજ્જત} \gu{છે}.
\end{quote}

\begin{quote}
\emph{gāmmā̃ tenā kuṭumbnī ghaṇī ijjat che.}
\end{quote}

\begin{quote}
His family is highly respected in the village.
\end{quote}

A feminine noun.

\begin{cousinweb}
\textbf{\gu{ઇજ્જત}} (\emph{ijjat}) is usually traced to Arabic \emph{ʿizzat}, \textquotedblleft{}might, honour\textquotedblright{}, through Persian.
\end{cousinweb}

\begin{grammarlens}[title={What you've built}]
One more word for giving reasons and opinions.
\end{grammarlens}

\subsection*{Your turn}
\begin{itemize}
\raggedright
  \item \emph{Say it:} \emph{ijjat}
  \item \emph{Say it:} \emph{ijjat}, once more
  \item \emph{Say it:} \emph{gāmmā̃ tenā kuṭumbnī ghaṇī ijjat che.} — His family is highly respected in the village.
  \item \emph{Recall:} say the Gujarati for cooperation, then the Gujarati for to keep up, then say \emph{ijjat} again
\end{itemize}

\begin{tcolorbox}[breakable,colback=teal!4,colframe=teal!35!black,title={Before you move on}]
What does \gu{ઇજ્જત} mean? (\textquotedblleft{}One's good name; honour\textquotedblright{}.) Say it once more.
\end{tcolorbox}
\section[\texorpdfstring{\gu{ગૌરવ} (gaurav)}{ગૌરવ (gaurav)}]{\gu{ગૌરવ} (gaurav) — dignity; pride in something worthy}

\label{lesson:GU-C359-gaurav}



\begin{quote}
Before the new one: say the Gujarati for to keep up, then the Gujarati for one's good name.
\end{quote}

\subsection*{You'll want to know: \gu{ગૌરવ}}
\textbf{\gu{ગૌરવ}} — \emph{gaurav} — \textquotedblleft{}dignity; pride in something worthy\textquotedblright{}.

\begin{quote}
\gu{અમને} \gu{અમારા} \gu{ગામનું} \gu{ગૌરવ} \gu{છે}.
\end{quote}

\begin{quote}
\emph{amne amārā gāmnũ gaurav che.}
\end{quote}

\begin{quote}
We are proud of our village.
\end{quote}

A neuter noun.

\begin{cousinweb}
\textbf{\gu{ગૌરવ}} (\emph{gaurav}) is a tatsama, borrowed from Sanskrit \emph{gaurava}, \textquotedblleft{}weight, importance, dignity\textquotedblright{} (Monier-Williams).
\end{cousinweb}

\begin{grammarlens}[title={What you've built}]
One more word for giving reasons and opinions.
\end{grammarlens}

\subsection*{Your turn}
\begin{itemize}
\raggedright
  \item \emph{Say it:} \emph{gaurav}
  \item \emph{Say it:} \emph{gaurav}, once more
  \item \emph{Say it:} \emph{amne amārā gāmnũ gaurav che.} — We are proud of our village.
  \item \emph{Recall:} say the Gujarati for to keep up, then the Gujarati for one's good name, then say \emph{gaurav} again
\end{itemize}

\begin{tcolorbox}[breakable,colback=teal!4,colframe=teal!35!black,title={Before you move on}]
What does \gu{ગૌરવ} mean? (\textquotedblleft{}Dignity; pride in something worthy\textquotedblright{}.) Say it once more.
\end{tcolorbox}
\section[\texorpdfstring{\gu{શરમ} (śaram)}{શરમ (śaram)}]{\gu{શરમ} (śaram) — shame}

\label{lesson:GU-C359-saram}



\begin{quote}
Before the new one: say the Gujarati for one's good name, then the Gujarati for dignity.
\end{quote}

\subsection*{You'll want to know: \gu{શરમ}}
\textbf{\gu{શરમ}} — \emph{śaram} — \textquotedblleft{}shame\textquotedblright{}.

\begin{quote}
\gu{ખોટું} \gu{બોલ્યા} \gu{પછી} \gu{તેને} \gu{શરમ} \gu{આવી}.
\end{quote}

\begin{quote}
\emph{khoṭũ bolyā pachī tene śaram āvī.}
\end{quote}

\begin{quote}
After telling a lie he felt ashamed.
\end{quote}

A feminine noun.

\begin{cousinweb}
\textbf{\gu{શરમ}} (\emph{śaram}) is usually traced to Persian \emph{sharm}, \textquotedblleft{}shame, modesty\textquotedblright{}. \textbf{\gu{શરમાવું}} (\emph{śarmāvũ}), \textquotedblleft{}to feel shy\textquotedblright{}, is built on it.
\end{cousinweb}

\begin{grammarlens}[title={What you've built}]
One more word for giving reasons and opinions.
\end{grammarlens}

\subsection*{Your turn}
\begin{itemize}
\raggedright
  \item \emph{Say it:} \emph{śaram}
  \item \emph{Say it:} \emph{śaram}, once more
  \item \emph{Say it:} \emph{khoṭũ bolyā pachī tene śaram āvī.} — After telling a lie he felt ashamed.
  \item \emph{Recall:} say the Gujarati for one's good name, then the Gujarati for dignity, then say \emph{śaram} again
\end{itemize}

\begin{tcolorbox}[breakable,colback=teal!4,colframe=teal!35!black,title={Before you move on}]
What does \gu{શરમ} mean? (\textquotedblleft{}Shame\textquotedblright{}.) Say it once more.
\end{tcolorbox}
\section[\texorpdfstring{\gu{અભિમાન} (abhimān)}{અભિમાન (abhimān)}]{\gu{અભિમાન} (abhimān) — conceit}

\label{lesson:GU-C359-abhiman}



\begin{quote}
Before the new one: say the Gujarati for dignity, then the Gujarati for shame.
\end{quote}

\subsection*{You'll want to know: \gu{અભિમાન}}
\textbf{\gu{અભિમાન}} — \emph{abhimān} — \textquotedblleft{}conceit\textquotedblright{}.

\begin{quote}
\gu{તેને} \gu{પોતાના} \gu{પૈસાનું} \gu{ઘણું} \gu{અભિમાન} \gu{છે}.
\end{quote}

\begin{quote}
\emph{tene potānā paisānũ ghaṇũ abhimān che.}
\end{quote}

\begin{quote}
He is very proud of his money.
\end{quote}

A neuter noun.

\begin{cousinweb}
\textbf{\gu{અભિમાન}} (\emph{abhimān}) is a tatsama, borrowed from Sanskrit \emph{abhimāna}, \textquotedblleft{}a high opinion of oneself, pride\textquotedblright{} (Monier-Williams).
\end{cousinweb}

\begin{grammarlens}[title={What you've built}]
One more word for giving reasons and opinions.
\end{grammarlens}

\subsection*{Your turn}
\begin{itemize}
\raggedright
  \item \emph{Say it:} \emph{abhimān}
  \item \emph{Say it:} \emph{abhimān}, once more
  \item \emph{Say it:} \emph{tene potānā paisānũ ghaṇũ abhimān che.} — He is very proud of his money.
  \item \emph{Recall:} say the Gujarati for dignity, then the Gujarati for shame, then say \emph{abhimān} again
\end{itemize}

\begin{tcolorbox}[breakable,colback=teal!4,colframe=teal!35!black,title={Before you move on}]
What does \gu{અભિમાન} mean? (\textquotedblleft{}Conceit\textquotedblright{}.) Say it once more.
\end{tcolorbox}
\section[\texorpdfstring{\gu{નમવું} (namvũ)}{નમવું (namvũ)}]{\gu{નમવું} (namvũ) — to bow}

\label{lesson:GU-C359-namvu}



\begin{quote}
Before the new one: say the Gujarati for shame, then the Gujarati for conceit.
\end{quote}

\subsection*{You'll want to know: \gu{નમવું}}
\textbf{\gu{નમવું}} — \emph{namvũ} — \textquotedblleft{}to bow\textquotedblright{}.

\begin{quote}
\gu{બાળકો} \gu{ગુરુ} \gu{આગળ} \gu{નમે} \gu{છે}.
\end{quote}

\begin{quote}
\emph{bāḷko guru āgaḷ name che.}
\end{quote}

\begin{quote}
The children bow before the teacher.
\end{quote}

Past (he …): \gu{તે} \gu{મંદિરમાં} \gu{નમ્યો}, \emph{te mandirmā̃ namyo}.

\begin{cousinweb}
\textbf{\gu{નમવું}} (\emph{namvũ}) is inherited from Sanskrit \emph{námati}, \textquotedblleft{}bends, bows\textquotedblright{}; CDIAL lists the Gujarati verb (CDIAL 6956).
\end{cousinweb}

\begin{grammarlens}[title={What you've built}]
One more word for giving reasons and opinions.
\end{grammarlens}

\subsection*{Your turn}
\begin{itemize}
\raggedright
  \item \emph{Say it:} \emph{namvũ}
  \item \emph{Say it:} \emph{namvũ}, once more
  \item \emph{Say it:} \emph{bāḷko guru āgaḷ name che.} — The children bow before the teacher.
  \item \emph{Recall:} say the Gujarati for shame, then the Gujarati for conceit, then say \emph{namvũ} again
\end{itemize}

\begin{tcolorbox}[breakable,colback=teal!4,colframe=teal!35!black,title={Before you move on}]
What does \gu{નમવું} mean? (\textquotedblleft{}To bow\textquotedblright{}.) Say it once more.
\end{tcolorbox}
